% =============================================================================
% HaluGuard: Final Report - ACL/ARR Style
% =============================================================================
\documentclass[11pt,a4paper]{article}

\usepackage[utf8]{inputenc}
\usepackage[T1]{fontenc}
\usepackage{times}
\usepackage{latexsym}
\usepackage{graphicx}
\usepackage{booktabs}
\usepackage{amsmath}
\usepackage{amssymb}
\usepackage{xcolor}
\usepackage{multirow}
\usepackage{url}
\usepackage{hyperref}
\usepackage{natbib}
\usepackage{capt-of}

\usepackage[a4paper,margin=2.5cm]{geometry}

\usepackage{multicol}
\setlength{\columnsep}{0.6cm}

\usepackage{titlesec}
\titleformat{\section}{\normalfont\large\bfseries}{\thesection}{1em}{}
\titleformat{\subsection}{\normalfont\normalsize\bfseries}{\thesubsection}{1em}{}
\titleformat{\subsubsection}{\normalfont\normalsize\itshape}{\thesubsubsection}{1em}{}
\titlespacing*{\section}{0pt}{2.5ex plus 1ex minus .2ex}{1.5ex plus .2ex}
\titlespacing*{\subsection}{0pt}{2ex plus 1ex minus .2ex}{1ex plus .2ex}

\setlength{\parindent}{0.4cm}
\setlength{\parskip}{0pt}

\allowdisplaybreaks
\sloppy
\setlength{\emergencystretch}{3em}

\renewenvironment{abstract}{%
    \begin{center}
    \textbf{Abstract}
    \end{center}
    \itshape
    \small
}{\vspace{0.5em}}

\makeatletter
\renewcommand{\maketitle}{%
    \begin{center}
    {\Large\bfseries \@title \par}
    \vspace{1.2em}
    {\normalsize \@author \par}
    \vspace{1.5em}
    \end{center}
}
\makeatother

\setcitestyle{authoryear,round,semicolon}

% =============================================================================
\title{HaluGuard: Execution-Grounded Contrastive Context Selection\\ for Hallucination-Free Repository-Level Code Generation}

\author{
  Sparsh Gupta\textsuperscript{1},
  Alan Liang\textsuperscript{1},
  Katherine Hancock\textsuperscript{1},
  Daniel Ho\textsuperscript{1},
  Freddie Liang\textsuperscript{1} \\[0.5em]
  \textsuperscript{1}\textit{University of Southern California} \\
  \texttt{\{sparshg, aqliang, kxhancoc, hodk, fhliang\}@usc.edu}
}

\date{}

\begin{document}
\maketitle

\begin{abstract}
Large Language Models for code generation suffer from \textit{hallucinations}---syntactically valid but semantically incorrect outputs that fabricate APIs, misuse identifiers, or violate dependencies. Existing retrieval-augmented approaches select context based on \textit{similarity} to the query, ignoring whether that context actually \textit{prevents} hallucinations. We propose \textbf{HaluGuard}, a framework that reframes context selection as a \textit{hallucination prevention problem}. We train seven Hallucination-Contrastive Context Selector (HCCS) architectures using listwise contrastive learning on RepoBench, supplemented by a rule-based Type-Specific Context Router and an Execution Feedback Loop (EFL). Our best model, DualEncoderDeep, achieves \textbf{29.0\%} exact match on code generation---a +9.0 percentage point gain over no-context and +4.0pp over BM25---and \textbf{85.6\%} Hard Acc@5 on RepoBench-R retrieval (+27.8pp over cosine similarity). A key finding is that \textit{embedding quality dominates architecture complexity}: a lightweight 800K-parameter dual-encoder atop frozen UniXcoder embeddings outperforms larger interaction-based models on downstream generation, suggesting the retrieval bottleneck lies in encoder representations rather than scorer capacity.
\end{abstract}

\begin{multicols}{2}

\section{Introduction}

Code generation using Large Language Models (LLMs) has achieved remarkable success, yet a critical failure mode persists: \textit{code hallucinations}---outputs that compile but produce incorrect behavior due to fabricated APIs, misused identifiers, or violated dependencies \citep{tian2025codehalu}. This problem intensifies in repository-level settings where models must reason about multi-file contexts with complex interdependencies.

Current approaches address this through retrieval-augmented generation (RAG), selecting context based on \textit{similarity} to the query using embeddings from models like CodeBERT \citep{liu2024graphcoder, phan2024repohyper}. However, this paradigm has a fundamental flaw: \textbf{similar code is not necessarily helpful code}. A function with high embedding similarity might share surface patterns but lack the precise type signatures, import statements, or API documentation needed to prevent hallucinations.

We propose a paradigm shift: instead of asking ``\textit{what context is similar to the query?}'', we ask ``\textit{what context prevents the model from hallucinating?}'' This reframing leads to \textbf{HaluGuard}, a framework with three novel components:

\textbf{(1) Hallucination-Contrastive Context Selector (HCCS):} We train seven lightweight scorer architectures to rank context by its \textit{hallucination prevention potential}, not similarity. Using listwise contrastive learning with gold-snippet supervision from RepoBench, HCCS learns to distinguish context that leads to correct code from context that induces hallucinations.

\textbf{(2) Type-Specific Context Routing:} Different hallucination types require different context. \textit{Resource} hallucinations (missing imports) need dependency information; \textit{naming} hallucinations need variable definitions; \textit{mapping} hallucinations need API documentation. We implement a lightweight rule-based boosting mechanism that upweights context chunks whose category matches the predicted hallucination type.

\textbf{(3) Execution Feedback Loop (EFL):} Inspired by recent work on execution-grounded RL for code \citep{gehring2024rlef}, we use execution results to iteratively refine context selection. When generated code fails, we analyze the failure mode, retrieve targeted context, and regenerate.

Our experiments on RepoBench show that HaluGuard's DualEncoderDeep model achieves 29.0\% exact match, outperforming cosine similarity by +2.0pp and BM25 by +4.0pp. On the RepoBench-R hard retrieval benchmark, HCCS models reach 85.6\% Acc@5 vs.\ 57.9\% for cosine---a 47.7\% relative improvement. On CrossCodeEval, HCCS improves exact match to 6.7\% (vs.\ 4.0\% no-context) and the v2 Type Router achieves the best identifier-level F1 (ID-F1=0.392), directly measuring reductions in naming hallucinations. Embedding quality (UniXcoder vs.\ CodeBERT) drives performance more than scorer architecture, and EFL is constrained by the SyntaxError-dominated failure distribution of fragment-level code benchmarks.

\section{Related Work}

\subsection{Code Hallucinations in LLMs}

\citet{tian2025codehalu} introduced CodeHalu, categorizing hallucinations into four types: \textit{mapping} (API misuse), \textit{naming} (identifier errors), \textit{resource} (missing imports), and \textit{logic} (incorrect algorithms). Their CodeHaluEval benchmark with 8,883 samples provides execution-based verification. \citet{zhang2025llmhallucinations} identified root causes and showed RAG reduces all hallucination types, motivating our focus on optimizing context selection.

\subsection{Context Selection for Code}

Existing approaches use similarity-based retrieval. \textbf{GraphCoder} \citep{liu2024graphcoder} and \textbf{RepoHyper} \citep{phan2024repohyper} build code graphs and retrieve via graph traversal. \textbf{DRaCo} \citep{cheng2024draco} uses dataflow analysis for precise retrieval. \textbf{CODEPROMPTZIP} \citep{he2025codepromptzip} compresses prompts while preserving critical tokens. \textbf{Hierarchical Context Pruning} \citep{zhang2024hcp} models repositories at the function level. None of these methods optimize for hallucination prevention directly---they all assume that similarity implies usefulness. Our work differs by explicitly training the selector to distinguish hallucination-inducing from hallucination-preventing context, measured through gold-snippet supervision from RepoBench.

\subsection{Contrastive Learning for Hallucination}

Contrastive decoding has emerged as a powerful technique for hallucination mitigation in vision-language models \citep{wang2024contrastive, leng2024mitigating}. These methods contrast outputs from ``good'' vs.\ ``bad'' conditions during inference. We adapt this principle to \textit{context selection during training}, learning to contrast hallucination-preventing vs.\ hallucination-inducing context through gold-snippet feedback---applying contrastive learning to retrieval rather than decoding.

\subsection{Execution-Guided Code Generation}

\citet{gehring2024rlef} showed that reinforcement learning with execution feedback dramatically improves code generation by grounding models in actual program behavior. \citet{le2022coderl} pioneered RL for code using test feedback. We extend this paradigm: execution results guide not just code generation, but also the retrieval of context that prevents failures.

\subsection{Repository-Level Code Generation}

Early benchmarks such as HumanEval focus on single-file tasks, failing to capture real-world multi-file complexity. \textbf{RepoBench} \citep{liu2023repobench} addresses this by introducing repository-level evaluation across retrieval, completion, and pipeline tasks. We adopt RepoBench as our primary benchmark as its cross-file retrieval setting directly stress-tests HaluGuard's hallucination-aware context selection under realistic dependency conditions. Unlike GraphCoder and RepoHyper, which require constructing explicit code graphs, our approach learns to rank existing context chunks without graph construction overhead.

\section{Dataset}

\textbf{RepoBench} \citep{liu2023repobench} is our primary benchmark. We use the \texttt{cross\_file\_first} Python split, containing \textbf{8,033 examples}, each with a \texttt{cropped\_code} query, a set of cross-file candidate snippets (mean $\approx$11 candidates), and a gold snippet index identifying the ground-truth context. We split examples 70/15/15 into train ($\sim$5,623), validation ($\sim$1,205), and test ($\sim$1,205) sets. For difficulty-stratified evaluation we define an \textit{easy} bucket (5--9 candidates, 359 examples) and a \textit{hard} bucket ($\geq$10 candidates, 522 examples).

\textbf{Training data:} We generate $\sim$80K contrastive triplets by pairing each query's gold snippet ($c^+$) against every other candidate ($c^-$). All queries and snippets are pre-encoded with two frozen backbone models---\textbf{CodeBERT} (microsoft/codebert-base) and \textbf{UniXcoder} (microsoft/unixcoder-base)---both producing 768-dimensional CLS-pooled embeddings. We embed queries with left-truncation (preserving recent code) and snippets with right-truncation (preserving snippet start), both at max 512 tokens. We evaluate two query views: \textit{full} (entire \texttt{cropped\_code}) and \textit{last3} (final 3 lines only).

\textbf{CrossCodeEval} \citep{ding2023crosscodeeval} serves as our secondary benchmark for ablating the Type Router and EFL. We use the Python split's \texttt{oracle\_bm25} configuration (150 sampled examples, seed=42), which provides exactly 5 pre-retrieved cross-file chunks per query. CrossCodeEval's real multi-file repositories and strict evaluation against gold completions make it a harder benchmark than RepoBench; exact match scores are consequently lower across all methods.

\section{Methods}

\subsection{Problem Formulation}

Given a coding query $q$ and repository context $\mathcal{C} = \{c_1, \ldots, c_n\}$, we select a subset $\mathcal{C}^* \subset \mathcal{C}$ that minimizes hallucination probability:
\begin{equation}
\mathcal{C}^* = \arg\min_{\mathcal{C}' \subset \mathcal{C}} P(\text{hallucination} \mid q, \mathcal{C}', \theta_{\text{LLM}})
\end{equation}
This differs fundamentally from similarity-based retrieval, which maximizes $\sum_{c \in \mathcal{C}'} \text{sim}(q, c)$. We operationalize hallucination prevention using the RepoBench gold snippet as a proxy: if a scorer ranks the gold snippet first, the LLM is more likely to receive the correct cross-file dependency.

\subsection{Hallucination-Contrastive Context Selector}

We train seven HCCS scorer architectures on pre-computed RepoBench embeddings. All architectures share a unified \texttt{score(query\_emb, chunk\_embs)} interface and are trained with masked listwise cross-entropy unless otherwise noted. The loss treats context selection as a ranking problem: the gold snippet $c^+$ is the correct class among all valid candidate chunks $\mathcal{C}$, with padding positions masked before the softmax:
\begin{equation}\label{eq:listwise}
\mathcal{L}_{\text{listwise}} = -\log \frac{\exp\!\big(s(q, c^+) / \tau\big)}{Z(q)}
\end{equation}
where $Z(q) = \sum_{c_i \in \mathcal{C}} \mathbb{1}[c_i \text{ valid}] \cdot \exp\!\big(s(q, c_i) / \tau\big)$
and $s(q, c) = \ell_2(E_{\phi}^{q}(q))^\top \cdot \ell_2(E_{\phi}^{c}(c))$
is the hallucination-prevention score via scaled cosine similarity, with $\tau$ a learnable temperature initialized at $e^{2.66} \approx 14.3$. The seven architectures are summarised in Table~\ref{tab:architectures}.

\begin{center}
\captionof{table}{HCCS scorer architectures evaluated.}
\label{tab:architectures}
\resizebox{\columnwidth}{!}{%
\begin{tabular}{clcrc}
\hline
\# & \textbf{Architecture} & \textbf{Loss} & \textbf{Params} & \textbf{Regularisation} \\
\hline
1 & DualEncoder        & Listwise CE & $\sim$200K  & LayerNorm, GELU, learned $\tau$ \\
2 & DualEncoderDeep    & Listwise CE & $\sim$800K  & Two-layer projection \\
3 & ListwiseMLP        & Listwise CE & $\sim$400K  & LayerNorm, GELU \\
4 & PairwiseMLP        & InfoNCE     & $\sim$400K  & ReLU + Sigmoid \\
5 & InteractionMLP     & Listwise CE & $\sim$1.1M  & Spectral norm, ESIM features \\
6 & BilinearScorer     & Listwise CE & $\sim$98K   & Factorised bilinear \\
7 & EnsembleScorer     & Listwise CE & tiny gate   & Frozen base + learned mixture \\
\hline
\end{tabular}}
\end{center}

\noindent\textbf{Training protocol.} We use AdamW ($\text{lr}=3\times10^{-4}$, weight decay $5\times10^{-3}$) with cosine annealing, batch size 256, and a maximum of 50 epochs with early stopping (patience 10). Regularisation includes dropout (0.3), label smoothing (0.1), and embedding dropout (0.15) applied to frozen pre-computed embeddings as an augmentation. We employ \textit{curriculum learning}: the first 5 warmup epochs use only random negatives, after which up to 4 hard negatives (cosine-similar non-gold chunks) are gradually introduced.

\subsection{Type-Specific Context Routing}

Different hallucination types require different remediation context: \textit{resource} $\rightarrow$ import/module snippets; \textit{naming} $\rightarrow$ class and function definitions; \textit{mapping} $\rightarrow$ typed function signatures; \textit{logic} $\rightarrow$ test and assertion patterns.

We implement a lightweight rule-based boosting mechanism using regex pattern analysis on \texttt{cropped\_code}. For each hallucination type $t \in \mathcal{T}$, a pre-emptive boost $\delta_t$ is predicted:
\begin{equation}
s'_i = \min\!\big(s_i + \delta_{t(i)},\; 1.0\big)
\end{equation}
where $s_i$ is the raw HCCS score for chunk $i$, $t(i)$ is the hallucination category assigned to chunk $i$ via \texttt{classify\_snippet}, and $\delta_{t(i)} \in [0.05, 0.20]$ is the additive boost for that category. In the post-failure (EFL) setting, boosts are derived from the actual Python exception raised during execution, mapped deterministically to a hallucination category via \texttt{ERROR\_TO\_CATEGORY}. Our v2 router uses mutually exclusive, priority-ordered categories (RESOURCE $\rightarrow$ MAPPING $\rightarrow$ LOGIC $\rightarrow$ NAMING) to prevent over-firing: naming patterns trigger only for external method calls, not arbitrary identifiers.

\subsection{Execution Feedback Loop}

Inspired by RLEF \citep{gehring2024rlef}, we iteratively refine context selection: (1) Generate code with current context; (2) Execute against tests; (3) If failure, classify error type (ImportError$\rightarrow$resource, NameError$\rightarrow$naming, etc.); (4) Retrieve targeted context for identified type; (5) Regenerate. This loop runs for up to $k=3$ iterations or until execution passes (10s timeout per attempt).

\subsection{Evaluation Metrics}

\textbf{Retrieval quality:} Accuracy@$k$ (Acc@$k$, $k \in \{1, 3, 5\}$) measures whether the gold snippet ranks in the top $k$; Mean Reciprocal Rank (MRR) rewards higher gold placement.

\textbf{Generation quality (RepoBench):} Exact Match (EM) for character-perfect code correctness; Edit Similarity (ES) for surface-level closeness; CodeBLEU (CB) for syntactic and semantic code quality.

\textbf{Generation quality (CrossCodeEval):} We additionally report BLEU-4 (4-gram corpus BLEU) and \textbf{ID-F1} (identifier-level F1 over extracted Python identifiers, excluding keywords and builtins). ID-F1 directly quantifies naming hallucination: a model that generates \texttt{client.fetch()} when the reference uses \texttt{client.get()} is penalised on the method name while receiving credit for the correctly predicted receiver.

\section{Experimental Results}

\subsection{Experimental Setup}

All retrieval experiments use pre-computed UniXcoder embeddings with the \textit{last3} query view (last 3 lines of context), which we found to outperform the full view on most architectures. Generation experiments use \textbf{Qwen2.5-Coder-7B} (4-bit NF4 quantization, A100-SXM4-80GB) with temperature 0.2, \texttt{max\_new\_tokens}=64, and top-5 retrieved chunks. RepoBench generation is evaluated on 200 held-out test examples; retrieval metrics are reported on the full test split ($\sim$1,205 examples) as well as the RepoBench-R protocol (881 examples with easy/hard bucketing). CrossCodeEval experiments use the same generator with TOP\_K=1 (HCCS selects the single best chunk from 5 oracle-BM25 candidates) on 150 sampled examples. Figure~\ref{fig:training_curves} shows validation MRR training curves across architectures.

\subsection{Retrieval Performance}

Table~\ref{tab:retrieval_test} reports retrieval metrics on our test split for all seven architectures plus the ensemble. InteractionMLP achieves the highest overall MRR (0.513) and Acc@5 (78.3\%), while the EnsembleScorer achieves the highest Acc@1 (32.1\%). Importantly, DualEncoder ($\sim$200K parameters) substantially outperforms BilinearScorer ($\sim$98K parameters) despite being a simpler model---demonstrating that the factorized bilinear underutilizes the 768-dim embedding space. All HCCS models considerably outperform the cosine similarity baseline, which achieves MRR $\approx 0.40$ and Acc@1 $\approx 18$\% (surpassed by all trained models).

\begin{center}
\captionof{table}{Architecture retrieval comparison on test split.}
\label{tab:retrieval_test}
\resizebox{\columnwidth}{!}{%
\begin{tabular}{lccccr}
\toprule
\textbf{Architecture} & \textbf{MRR} & \textbf{Acc@1} & \textbf{Acc@3} & \textbf{Acc@5} & \textbf{Params} \\
\midrule
InteractionMLP   & \textbf{0.513} & 31.0 & \textbf{63.7} & \textbf{78.3} & 1.1M \\
EnsembleScorer   & 0.512 & \textbf{32.1} & 62.4 & 76.3 & small \\
DualEncoderDeep  & 0.489 & 29.6 & 59.9 & 74.9 & 800K \\
DualEncoder      & 0.477 & 28.2 & 58.8 & 72.2 & 200K \\
ListwiseMLP      & 0.451 & 24.2 & 57.5 & 72.5 & 400K \\
PairwiseMLP      & 0.406 & 20.2 & 51.1 & 67.1 & 400K \\
BilinearScorer   & 0.382 & 16.8 & 48.6 & 65.2 & 98K \\
\bottomrule
\end{tabular}}
\end{center}

Table~\ref{tab:repobench_r} shows RepoBench-R results stratified by difficulty. DualEncoderDeep achieves the highest absolute accuracy on both easy (68.2\% Acc@1) and hard (85.6\% Acc@5) splits. The Ensemble matches or exceeds DualEncoderDeep on easy Acc@3 (83.8\% vs.\ 82.5\%) owing to diversity in the learned mixture. Compared to the cosine baseline (UniXcoder Last-3), trained HCCS models improve Hard Acc@1 by an absolute +45.8pp (DualEncoderDeep 64.4\% vs.\ cosine 18.6\%)---demonstrating that learned hallucination-aware scoring is especially advantageous when the candidate pool is large and similarity-based signals become less discriminative.

\begin{center}
\captionof{table}{RepoBench-R easy/hard accuracy (\%). Easy: 5--9 candidates (359 ex.); Hard: $\geq$10 candidates (522 ex.).}
\label{tab:repobench_r}
\resizebox{\columnwidth}{!}{%
\begin{tabular}{lccccc}
\toprule
\textbf{Method} & \multicolumn{2}{c}{\textbf{Easy}} & \multicolumn{3}{c}{\textbf{Hard}} \\
\cmidrule(lr){2-3}\cmidrule(lr){4-6}
 & @1 & @3 & @1 & @3 & @5 \\
\midrule
DualEncoderDeep  & \textbf{68.2} & 82.5 & \textbf{64.4} & \textbf{79.5} & \textbf{85.6} \\
EnsembleScorer   & 64.4 & \textbf{83.8} & 64.8 & 79.7 & 84.9 \\
DualEncoder      & 51.3 & 78.0 & 46.9 & 69.7 & 78.2 \\
ListwiseMLP      & 41.0 & 74.9 & 29.5 & 53.1 & 68.4 \\
InteractionMLP   & 38.2 & 73.5 & 30.1 & 56.7 & 70.3 \\
Cosine (UniX L3) & 27.9 & 61.3 & 18.6 & 41.0 & 57.9 \\
Random           & 15.3 & 45.6 & ~6.3 & 19.0 & 32.0 \\
\bottomrule
\end{tabular}}
\end{center}

\noindent A notable inversion appears between the two tables: InteractionMLP leads on MRR in the full test split yet falls behind DualEncoderDeep on the RepoBench-R protocol. We attribute this to the latter's difficulty-weighted evaluation---InteractionMLP's ESIM interaction features improve fine-grained ranking across all candidates (boosting MRR), but the dual-encoder's cleaner cosine geometry is more robust when the candidate pool is large.

\begin{center}
\includegraphics[width=\columnwidth]{data/results/training_curves.png}
\captionof{figure}{Validation MRR over 50 training epochs for six HCCS architectures. DualEncoderDeep (orange) separates clearly from the field after epoch~15, converging to MRR~$\approx$0.42. Most architectures plateau near 0.39--0.40; InteractionMLP (purple) continues improving through epoch~30.}
\label{fig:training_curves}
\end{center}

\subsection{RepoBench: Generation Quality}

Table~\ref{tab:generation} shows end-to-end code generation results on 200 RepoBench test examples with Qwen2.5-Coder-7B. DualEncoderDeep achieves the highest exact match (29.0\%) and strong CodeBLEU (0.562), followed by ListwiseMLP (28.5\% EM). The performance order on generation does not perfectly mirror the retrieval ranking, highlighting that retrieval quality and generation quality are related but not identical objectives.

\begin{center}
\captionof{table}{RepoBench generation results (n=200). Sorted by EM.}
\label{tab:generation}
\resizebox{\columnwidth}{!}{%
\begin{tabular}{lccc}
\toprule
\textbf{Method} & \textbf{EM} & \textbf{ES} & \textbf{CB} \\
\midrule
\multicolumn{4}{l}{\textit{HCCS Models}} \\
DualEncoderDeep  & \textbf{0.290} & \textbf{0.646} & 0.562 \\
ListwiseMLP      & 0.285 & 0.634 & 0.556 \\
EnsembleScorer   & 0.275 & 0.625 & 0.550 \\
DualEncoder      & 0.270 & 0.632 & 0.554 \\
InteractionMLP   & 0.270 & 0.627 & 0.555 \\
BilinearScorer   & 0.270 & 0.632 & 0.554 \\
PairwiseMLP      & 0.265 & 0.623 & 0.547 \\
\midrule
\multicolumn{4}{l}{\textit{Baselines}} \\
Cosine UniX (L3) & 0.270 & 0.628 & 0.554 \\
Full Context     & 0.265 & 0.645 & \textbf{0.563} \\
BM25             & 0.250 & 0.627 & 0.543 \\
Gold Only        & 0.250 & 0.614 & 0.537 \\
No Context       & 0.200 & 0.576 & 0.500 \\
\bottomrule
\end{tabular}}
\end{center}

Several findings merit attention. \textbf{BM25} is a standard lexical RAG baseline: it retrieves context via keyword overlap, passes it to the LLM, and the LLM generates code conditioned on that context---the same pipeline used by all methods here, differing only in how context is selected. Despite its widespread use, BM25 achieves only 25.0\% EM, which every HCCS model surpasses. \textbf{Gold Only} (oracle: always provide the ground-truth snippet) also achieves only 25.0\% EM---lower than most HCCS models. The oracle \textit{single snippet} alone is insufficient; HCCS's top-5 selection provides richer cross-file information even when it includes non-gold snippets. \textbf{Full Context} achieves only 26.5\% EM, confirming that more context is not always better---excess irrelevant snippets dilute the generative prompt. The gap between No Context (20.0\%) and DualEncoderDeep (29.0\%) represents a +9.0pp absolute improvement from context selection alone.

\subsection{CrossCodeEval: Type Router and EFL Ablation}

To isolate the contributions of the Type-Specific Context Router and EFL, we evaluate six conditions on the CrossCodeEval Python benchmark \citep{ding2023crosscodeeval} (n=150 examples, seed=42). We use the \texttt{oracle\_bm25} split which provides exactly 5 pre-retrieved cross-file chunks per query; HCCS selects the single top-ranked chunk (TOP\_K=1) from these 5, making the scorer's discrimination the sole driver of context quality. We additionally report \textbf{BLEU-4} (4-gram corpus BLEU) and \textbf{ID-F1} (identifier-level F1), which directly measures naming hallucination: if the model generates \texttt{client.fetch()} instead of \texttt{client.get()}, ID-F1 penalises the incorrect method name while crediting the correctly predicted receiver object.

\begin{center}
\captionof{table}{CrossCodeEval ablation (n=150). ListwiseMLP scorer, TOP\_K=1. EFL pass@1 reported for condition~6.}
\label{tab:cceval}
\resizebox{\columnwidth}{!}{%
\begin{tabular}{lcccc}
\toprule
\textbf{Condition} & \textbf{EM} & \textbf{ES} & \textbf{BLEU-4} & \textbf{ID-F1} \\
\midrule
(1) No context                         & 0.040 & 0.527 & 0.120 & 0.363 \\
(2) BM25 baseline                      & 0.053 & \textbf{0.537} & 0.095 & 0.390 \\
(3) HCCS, no router                    & \textbf{0.067} & 0.533 & 0.119 & 0.383 \\
(4) HCCS + v1 router                   & \textbf{0.067} & 0.528 & \textbf{0.127} & 0.378 \\
(5) HCCS + v2 router                   & 0.053 & 0.530 & 0.110 & \textbf{0.392} \\
(6) HCCS + v2 + EFL\textsuperscript{$\dagger$} & \textbf{0.067} & 0.531 & 0.121 & 0.377 \\
\bottomrule
\multicolumn{5}{l}{\footnotesize $\dagger$ EFL pass@1 = 0.007 (1/150 examples); see discussion.} \\
\end{tabular}}
\end{center}

\noindent\textbf{Router analysis.} The v1 router fires on \textbf{100\%} of CrossCodeEval queries: the NAMING and RESOURCE patterns both trigger for every example. Despite this over-firing, condition~(4) matches condition~(3) on EM (0.067 = 0.067) and achieves higher BLEU-4 (0.127 vs.\ 0.119), suggesting that uniform boosting happens to reinforce the scorer's ordering on this split. However, v1 degrades ID-F1 (0.378 vs.\ 0.383), indicating that naming-specific correctness suffers when all snippets receive the same boost regardless of content type.

The v2 router (mutually exclusive, priority-ordered) assigns categories to 96.7\% of queries: RESOURCE (32.0\%), MAPPING (34.7\%), NAMING (17.3\%), LOGIC (12.7\%), with 3.3\% receiving no boost. Condition~(5) drops EM to 0.053 ($-$1.4pp vs.\ no-router) yet achieves the \textbf{best ID-F1} (0.392, $+$0.9pp). Per-category EM delta reveals the source of the trade-off: MAPPING queries benefit ($+$1.9pp EM), while RESOURCE ($-$4.2pp) and LOGIC ($-$5.3pp) categories are hurt---the additive boost disrupts the scorer's learned geometry for non-naming categories. These results motivate replacing rule-based boosting with a learned router that adjusts boost magnitudes per category.

\noindent\textbf{EFL analysis.} Condition~(6) adds the Execution Feedback Loop. EFL achieves a pass@1 of \textbf{0.007} (1 of 150 examples): 149 of 150 queries exhaust all 3 retry iterations without passing. Error analysis over 448 execution attempts reveals a severe benchmark mismatch: \textbf{81.2\% SyntaxError}, 13.2\% ModuleNotFoundError, 3.3\% IndentationError, 2.0\% ImportError---and \textbf{0\%} NameError or AttributeError. CrossCodeEval provides code \textit{fragments} rather than fully executable files; structurally malformed outputs produce SyntaxErrors that no amount of context substitution can resolve. EFL's primary design targets---NameError and AttributeError---are entirely absent from the failure distribution. Only ImportError (2.0\%) falls within EFL's actionable scope. The mechanism is architecturally sound, but the benchmark's fragment-level execution environment does not exercise the failure modes EFL was designed to address; evaluation on full-repository benchmarks such as SWE-bench would expose the naming and import errors that drive EFL's value.

\section{Conclusion and Future Work}

We presented HaluGuard, a framework that reframes repository-level context selection as a hallucination prevention problem. By training seven HCCS scorer architectures via listwise contrastive learning on RepoBench, we demonstrated consistent improvements over similarity-based retrieval. Our DualEncoderDeep model achieves 29.0\% exact match on code generation (+9.0pp over no context, +4.0pp over BM25) and 85.6\% Hard Acc@5 on retrieval (+27.8pp over cosine similarity). On CrossCodeEval, the v2 Type Router achieves the highest ID-F1 (0.392), confirming that category-specific boosting reduces naming hallucinations even when it trades off some exact match.

The central finding of this work is that \textbf{embedding quality is the dominant bottleneck, not scorer complexity}: UniXcoder-based encoders outperform CodeBERT-based ones by $\sim$5 MRR points, whereas architecture variations within the same encoder yield differences $<$1 MRR point. A 200K-parameter DualEncoder matches or beats a 1.1M-parameter InteractionMLP on generation despite the latter's superior retrieval MRR, suggesting that downstream generation quality is less sensitive to top-$k$ reranking precision than to the quality of the retrieved context set as a whole. The Type Router's architecture sensitivity---helps DualEncoder, hurts DualEncoderDeep on RepoBench---further supports this: additive boosting disrupts deep projection geometry more than shallow cosine scoring.

\textbf{Future work} should pursue five directions: (1) \textit{Encoder fine-tuning via adapter layers}---lightweight adapters on UniXcoder could directly optimise for hallucination-prevention rather than general code similarity; (2) \textit{Execution-verified training labels}---replacing RepoBench gold-snippet supervision with execution pass/fail signals as in RLEF \citep{gehring2024rlef} would provide a more direct training signal; (3) \textit{Learned type routing}---replacing the rule-based additive booster with a lightweight learned router would enable continuous score adjustment without disrupting the scorer's geometry; (4) \textit{Broader benchmark coverage}---evaluation on SWE-bench and HumanEval-Pack with fully executable repositories would better characterise EFL's utility; (5) \textit{Multi-language evaluation}---extending to Java and TypeScript within RepoBench to assess cross-language generality.

\section{Division of Labor}

\noindent\textbf{Sparsh Gupta:} Worked on training, regularization, and designing of three HCCS architectures; designed the Type-Specific Context Router and EFL architecture; conducted evaluation of the retrieval pipeline using RepoBench and the generation pipeline using CrossCodeEval.

\noindent\textbf{Alan Liang:} Researched alternative training objectives, loss schedules, and regularization; ran ablation studies on generalization; worked on the Type-Specific Context Router.

\noindent\textbf{Katherine Hancock:} Generated all CodeBERT and UniXcoder embeddings; trained HCCS variants; wrote the initial draft of this report; worked on the Type-Specific Context Router.

\noindent\textbf{Daniel Ho:} Explored MLP scorer improvements including alternative features, score distribution logging, and warmup/scheduler tuning; researched BigCode for pass@k; worked on authoring the final report.

\noindent\textbf{Freddie Liang:} Implemented four scorer architectures across both encoders (8 experiments); ran experiments on other relevant generation benchmarks.

\bibliographystyle{plainnat}
\begin{thebibliography}{13}

\bibitem[Cheng et al.(2024)]{cheng2024draco}
Cheng, W., et al. (2024).
\newblock DRaCo: Dataflow-Guided Retrieval Augmentation for Repository-Level Code Completion.
\newblock \textit{ACL Anthology}.

\bibitem[Ding et al.(2023)]{ding2023crosscodeeval}
Ding, Y., et al. (2023).
\newblock CrossCodeEval: A Diverse and Multilingual Benchmark for Cross-File Code Completion.
\newblock \textit{arXiv:2310.11248}.

\bibitem[Gehring et al.(2024)]{gehring2024rlef}
Gehring, J., et al. (2024).
\newblock RLEF: Grounding Code LLMs in Execution Feedback with Reinforcement Learning.
\newblock \textit{ICML 2025}.

\bibitem[He et al.(2025)]{he2025codepromptzip}
He, P., Wang, S., and Chen, T.-H. (2025).
\newblock CODEPROMPTZIP: Code-specific Prompt Compression for RAG.
\newblock \textit{arXiv:2502.14925}.

\bibitem[Le et al.(2022)]{le2022coderl}
Le, H., et al. (2022).
\newblock CodeRL: Mastering Code Generation through Pretrained Models and Deep RL.
\newblock \textit{NeurIPS}.

\bibitem[Leng et al.(2024)]{leng2024mitigating}
Leng, S., et al. (2024).
\newblock Mitigating Object Hallucinations via Contrastive Decoding.
\newblock \textit{CVPR}.

\bibitem[Liu et al.(2023)]{liu2023repobench}
Liu, T., et al. (2023).
\newblock RepoBench: Benchmarking Repository-Level Code Auto-Completion Systems.
\newblock \textit{arXiv:2306.03091}.

\bibitem[Liu et al.(2024)]{liu2024graphcoder}
Liu, W., et al. (2024).
\newblock GraphCoder: Enhancing Repository-Level Code Completion via Code Context Graph-based Retrieval.
\newblock \textit{arXiv preprint}.

\bibitem[Phan et al.(2024)]{phan2024repohyper}
Phan, H.~N., et al. (2024).
\newblock RepoHyper: Search-Expand-Refine on Semantic Graphs for Repository-Level Code Completion.
\newblock \textit{arXiv:2403.06095}.

\bibitem[Tian et al.(2025)]{tian2025codehalu}
Tian, Y., et al. (2025).
\newblock CodeHalu: Investigating Code Hallucinations in LLMs via Execution-based Verification.
\newblock \textit{AAAI}.

\bibitem[Wang et al.(2024)]{wang2024contrastive}
Wang, X., et al. (2024).
\newblock Mitigating Hallucinations in LVLMs with Instruction Contrastive Decoding.
\newblock \textit{ACL Findings}.

\bibitem[Zhang et al.(2024)]{zhang2024hcp}
Zhang, L., et al. (2024).
\newblock Hierarchical Context Pruning: Optimizing Real-World Code Completion.
\newblock \textit{arXiv:2406.18294}.

\bibitem[Zhang et al.(2025)]{zhang2025llmhallucinations}
Zhang, Z., et al. (2025).
\newblock LLM Hallucinations in Practical Code Generation: Phenomena, Mechanism, and Mitigation.
\newblock \textit{arXiv:2409.20550}.

\end{thebibliography}

\end{multicols}

\end{document}
